\documentclass[11pt]{article}
\usepackage[a4paper,margin=28mm]{geometry}
\usepackage{amsmath,amssymb,amsthm,mathtools}
\usepackage[T1]{fontenc}
\usepackage{lmodern,microtype}
\usepackage[colorlinks=true,linkcolor=blue,citecolor=blue,urlcolor=blue]{hyperref}
\usepackage{enumitem}
\setlist{nosep}
\newtheorem{theorem}{Theorem}[section]
\newtheorem{proposition}[theorem]{Proposition}
\newtheorem{lemma}[theorem]{Lemma}
\newtheorem{corollary}[theorem]{Corollary}
\theoremstyle{remark}\newtheorem{remark}[theorem]{Remark}
\newcommand{\Li}{\operatorname{Li}}
\newcommand{\Prob}{\mathbb P}
\newcommand{\E}{\mathbb E}
\newcommand{\one}{\mathbf 1}
\newcommand{\cP}{\mathcal P}
\newcommand{\cA}{\mathcal A}
\newcommand{\cQ}{\mathcal Q}
\newcommand{\eps}{\varepsilon}
\newcommand{\FI}{F_{\mathrm I}}
\newcommand{\FII}{F_{\mathrm {II}}}
\newcommand{\cB}{\mathcal B}
\title{Complete local counts at prescribed initial positions\newline
in Erd\H{o}s--Straus decompositions}
\author{}
\date{October 30, 2026}
\begin{document}
\maketitle
\begin{abstract}
For primes $p\equiv1\pmod4$, we count all unordered Type~I and Type~II
Erd\H{o}s--Straus decompositions at prescribed initial smallest
denominators. Put $\delta=2^{-20}$ and $\kappa=(\log3)/8-\delta$.
Outside one set of $O(X/(\log X)^{1+\delta})$ primes up to $X$,
uniformly for $a\equiv3\pmod4$, $3\le a\le(\log p)^\delta$, and
$x=(p+a)/4$, the two counts are
\[
 \FI(p,x)=\frac{\tau(x^2)}{\varphi(a)}
      \bigl(1+O((\log p)^{-\kappa})\bigr),\qquad
 \FII(p,x)=\frac{\tau(x^2)}{2\varphi(a)}
      \bigl(1+O((\log p)^{-\kappa})\bigr).
\]
On a possibly smaller common good set with the same exceptional-set
order, each count is bounded below by a constant multiple of
$(\log p)^{\theta\log3}/\varphi(a)$ for any fixed
$0<\theta<15/16$. The argument combines established divisor
parametrizations, finite-group character estimates, and uniform
distribution estimates for shifted primes. We give both a direct
support-compressed moment proof and a uniform application of Elliott's
Bernoulli model.
\end{abstract}
\noindent\textbf{Keywords.} Egyptian fractions; prescribed denominators;
local multiplicities; shifted primes; divisor distributions;
Bombieri--Vinogradov theorem.

\section{Introduction and main results}
The Erd\H{o}s--Straus equation is
\begin{equation}\label{eq:es}
 \frac4p=\frac1x+\frac1y+\frac1z,
 \qquad x,y,z\in\mathbb Z_{>0}.
\end{equation}
We take $p\equiv1\pmod4$ prime and prescribe initial positions
$x=(p+a)/4$ with $a=3+4j$. Solutions are counted without regard to the
order of their denominators. Write $\FI(p,x)$, respectively
$\FII(p,x)$, for the number of solutions with strictly smallest
denominator $x$ and precisely one, respectively two, denominators
divisible by $p$. All logarithms are natural, and $\tau(n)$ denotes the
number of positive divisors of $n$.

For $3\le a<p$, $a\equiv3\pmod4$, we have
\[
 (a,p)=(a,x)=1,\qquad x<p/2,\qquad p\nmid x.
\]
Reduction of the cleared equation modulo $p$ shows that any solution
at this $x$ has at least one of its other denominators divisible by $p$.
Thus the two types exhaust the complete local count at this position.

\begin{theorem}[Complete local counts]\label{thm:main}
Set
\[
\begin{gathered}
 \delta=2^{-20},\qquad \kappa=\frac{\log3}{8}-\delta>0,\\
 J(p)=\left\lfloor\frac{(\log p)^\delta+1}{4}\right\rfloor,
 \qquad a_j=3+4j,\quad x_j=(p+a_j)/4.
\end{gathered}
\]
There is a set $\mathcal E$ of primes $p\equiv1\pmod4$ such that
\begin{equation}\label{eq:global-exception}
 \#\{p\le X:p\in\mathcal E\}\ll\frac{X}{(\log X)^{1+\delta}}
 \qquad(X\to\infty),
\end{equation}
and, for sufficiently large $p\notin\mathcal E$, simultaneously at
all positions $0\le j<J(p)$,
\begin{align}
 \FI(p,x_j)
 &=\frac{\tau(x_j^2)}{\varphi(a_j)}
       \left(1+O((\log p)^{-\kappa})\right),\label{eq:main-I}\\
 \FII(p,x_j)
 &=\frac{\tau(x_j^2)}{2\varphi(a_j)}
       \left(1+O((\log p)^{-\kappa})\right).\label{eq:main-II}
\end{align}
The error constants are uniform in $j$. In particular,
\begin{equation}\label{eq:ratio-total}
\begin{aligned}
 \frac{\FI(p,x_j)}{\FII(p,x_j)}
 &=2+O((\log p)^{-\kappa}),\\
 \FI(p,x_j)+\FII(p,x_j)
 &=\frac{3\tau(x_j^2)}{2\varphi(a_j)}
       \left(1+O((\log p)^{-\kappa})\right).
\end{aligned}
\end{equation}
Solutions at different positions $j$ are distinct.
\end{theorem}

\begin{corollary}[Quantitative multiplicity]\label{cor:many}
Fix $0<\theta<15/16$ and put $\gamma_\theta=\theta\log3$.
There is an exceptional set $\mathcal E_\theta\supseteq\mathcal E$ with
\[
 \#\{p\le X:p\in\mathcal E_\theta\}
 \ll_\theta X/(\log X)^{1+\delta}
\]
such that, for sufficiently large $p\notin\mathcal E_\theta$, all
positions $0\le j<J(p)$ satisfy Theorem~\ref{thm:main} and
\begin{equation}\label{eq:many}
 \FI(p,x_j)\ge\frac{(\log p)^{\gamma_\theta}}{4\varphi(a_j)},
 \qquad
 \FII(p,x_j)\ge\frac{(\log p)^{\gamma_\theta}}{8\varphi(a_j)}.
\end{equation}
In particular, one may take
$\gamma_*=\gamma_{7/8}=7\log3/8=0.96128575\ldots$.
\end{corollary}
The enlarged exceptional set comes from one additional total-factor
statistic per shift. Its order is unchanged, but it need not be the
same set point by point. The range $\theta<15/16$ is sufficient for the
chosen proof parameters; no endpoint or optimality assertion is made.

\begin{corollary}[Fixed initial catalogue]\label{cor:fixed}
For each fixed integer $J_0\ge1$, the local asymptotics
\eqref{eq:main-I}--\eqref{eq:main-II} hold simultaneously for
$0\le j<J_0$ outside an exceptional set of
$O_{J_0}(X/(\log X)^{1+\delta})$ primes $p\le X$, $p\equiv1\pmod4$.
For each additional fixed $0<\theta<15/16$, the bounds \eqref{eq:many}
hold for these positions outside a possibly enlarged exceptional set
of size $O_{J_0,\theta}(X/(\log X)^{1+\delta})$.
\end{corollary}

The constants are deliberately unoptimized. The catalogue becomes
nonempty precisely when
\begin{equation}\label{eq:nonempty}
 J(p)\ge1\quad\Longleftrightarrow\quad p\ge\exp(3^{2^{20}}).
\end{equation}
The results are unconditional and asymptotic, with a further
ineffective onset inherited from Siegel--Walfisz. The exceptional set
is defined by sufficient factor-distribution tests. The theorem leaves
the individual exceptional primes untreated and gives no complexity
bound for constructing the decompositions.

\subsection{Earlier results and the quantitative application}
Bradford's Propositions~1--4 \cite{Bradford} give complete
fixed-denominator divisor criteria for both types. We use these
criteria in a normalized form and include their inverse maps to fix
the counting conventions. Elsholtz and Tao
\cite[Theorem~1.8 and Section~6]{ET} already obtain many solutions at
the first position $x=(p+3)/4$ for a relative-density-one set of
primes. Their fixed-shift additive-function theorem
\cite[Appendix, Lemma~A.3]{ET}, together with the complete divisor
criteria and character orthogonality, also gives complete local
asymptotics on any fixed finite catalogue. It yields multiplicity
exponents arbitrarily close to $\log3$ on such a catalogue.
Appendix~\ref{app:fixed} records this deduction explicitly.

Hall's work on the divisors of $p-1$ \cite{Hall} and the per-class
divisor estimates in the proof of Bettin et al., Lemma~4.9
\cite{Bettin}, provide further fixed-shift precedents. The latter
also give a one-sided Type~II lower bound of exponent
$\log2-\eps$ on every fixed finite catalogue; the substitution is
given in Appendix~\ref{app:fixed}. A diagonal choice from fixed-catalogue
results allows some unspecified slowly growing catalogue. Thus the
quantitative content considered here is the prescribed logarithmic-power
length, a uniform relative error for both complete counts, and one
common exceptional-set estimate.

Other nearby results concern different counting questions.
Elsholtz and Planitzer \cite{EP} obtain lower bounds for total solution
counts on density-one sets of integers and on infinite sets of primes.
Pomerance and Weingartner \cite{PW} bound exceptional sets for
existence of a decomposition with a varying numerator. Dahan
\cite[Section~4.3]{Dahan} allows a growing search depth in a Type~II
witness family. Complete parametrizations are also studied in
Bello--Hern\'andez, Benito, and Fern\'andez \cite{Bello}, while
Chamberland \cite{Chamberland} gives a Type~II prime-representation
criterion. These conclusions concern total counts, existence, or
search depth; the present theorems count each prescribed position.

Our proof uses standard distribution results for primes. A direct
moment expansion groups terms by distinct squarefree supports.
The normalized coefficient of each support is polylogarithmic,
and its progression modulus is at most $4N^{1/4}$, so ordinary
Bombieri--Vinogradov controls the accumulated error.
Alternatively, Elliott's established Bernoulli model, in the form
recalled by Gorodetsky \cite[Theorem~1.1]{Gorodetsky}, gives the same
transfer after a uniform remainder calculation;
Appendix~\ref{app:elliott} supplies that calculation. Ford's
Poisson model for fixed shifts \cite{Ford} is another relevant
distributional framework. The application here combines these
analytic inputs with complete prime-power divisor boxes and the
exact Type~I/Type~II multiplicities.

\section{Complete local counts and character mixing}\label{sec:dictionary}
\subsection{The square-divisor dictionary}
For this subsection, assume $p\equiv1\pmod4$ is prime,
$a\equiv3\pmod4$, $3\le a<p$, and $x=(p+a)/4$.
In particular $(a,x)=(a,p)=1$ and $x<p/2$.

For $b\in(\mathbb Z/a\mathbb Z)^\times$, define
\begin{equation}\label{eq:C}
 C_{a,x}(b)=\#\{d\mid x^2:d\,x^{-1}\equiv b\pmod a\},
 \qquad M(x)=\tau(x^2).
\end{equation}
The inverse in \eqref{eq:C} is taken modulo $a$.

\begin{lemma}[Divisors and primitive ratios]\label{lem:dictionary}
The divisors $d\mid x^2$ correspond bijectively to ordered pairs
$(u,v)$ with $(u,v)=1$ and $uv\mid x$. The correspondence is
\[
 g=(d,x),\quad u=d/g,\quad v=x/g,\quad D=g^2/d,
 \qquad x=Duv,\quad d=Du^2.
\]
In particular, $D$ is a positive integer and $d/x=u/v$.
If $x=\prod_{q\mid x}q^{e_q}$ and
$d=\prod_{q\mid x}q^{j_q}$, the ratio exponent at $q$ is
$j_q-e_q\in\{-e_q,\ldots,e_q\}$.
\end{lemma}
\begin{proof}
At each prime, the exponents of $u,v,D$ are respectively
$\max(j_q-e_q,0)$, $\max(e_q-j_q,0)$, and
$2\min(j_q,e_q)-j_q$. The last is nonnegative because
$0\le j_q\le2e_q$. These formulas give all the claimed identities.
Conversely, a coprime pair with $uv\mid x$ determines $D=x/(uv)$ and
$d=Du^2$, which divides $x^2$ and reverses the construction.
\end{proof}

\begin{proposition}[Exact counts]\label{prop:exact}
In the stated range,
\begin{equation}\label{eq:exact}
 \FI(p,x)=C_{a,x}(-p^{-1}),\qquad
 \FII(p,x)=\frac12 C_{a,x}(-1).
\end{equation}
Both identities count all solutions with strictly smallest denominator $x$.
\end{proposition}
\begin{proof}
For Type~I, the condition $d x^{-1}\equiv-p^{-1}\pmod a$ is equivalent
to $pd+x\equiv0\pmod a$, or to $4d+1\equiv0\pmod a$.
Since $dd'=x^2$ for $d'=x^2/d$, it also implies
$d'+px\equiv0\pmod a$. Set
\begin{equation}\label{eq:typeIdiv}
 y=\frac{px+x^2/d}{a},\qquad z=\frac{p(x+pd)}{a}.
\end{equation}
Equivalently, with Lemma~\ref{lem:dictionary},
\begin{equation}\label{eq:typeIpair}
 c_1=\frac{pu+v}{a},\qquad
 (x,y,z)=(Duv,Dv c_1,pDu c_1).
\end{equation}
Indeed,
\[
 \frac1y+\frac1z=\frac{pu+v}{pDuvc_1}=\frac a{px},
 \qquad \frac1x+\frac a{px}=\frac4p.
\]
Moreover $p\nmid Duv$ and
$ac_1\equiv v\not\equiv0\pmod p$, so only $z$ is divisible by $p$.
We have
\[
 \frac yx=\frac{pu+v}{au}>\frac pa>1,\qquad
 \frac zy=\frac{pu}{v}\ge\frac p x>2.
\]
Thus $x<y<z$. The ordered pair $(u,v)$ is recovered as the reduced ratio
$z/(py)$, or $d$ is recovered from $d=(az/p-x)/p$.
There is no division by two.

For completeness, take any Type~I solution and name the non-$p$ second
denominator $y$ and the $p$-divisible one $pV$.
The reciprocal identity is
$a y V=pxV+xy$. It implies $aV-x>0$ and, after reduction modulo $p$,
$p\mid aV-x$. Write $aV-x=pd$ with $d>0$. Substitution gives
\[
 (ay-px)d=x^2.
\]
Hence $d\mid x^2$, $pd+x\equiv0\pmod a$, and the solution is
\eqref{eq:typeIdiv}. The inequalities already proved also show that its
$p$-divisible denominator is the larger one.

For Type~II write the two other denominators as $pY,pZ$. Then
\[
 \frac ax=\frac1Y+\frac1Z,\qquad
 (aY-x)(aZ-x)=x^2.
\]
Both factors $aY-x$ and $aZ-x$ are positive by the reciprocal
identity. This gives all divisors $d\mid x^2$ with $d+x\equiv0\pmod a$, by
\begin{equation}\label{eq:typeIIdiv}
 \{y,z\}=
 \left\{\frac{p(x+d)}a,\frac{p(x+x^2/d)}a\right\}.
\end{equation}
The two denominators are positive multiples of $p$, so they exceed
$x<p/2$. Swapping them sends $d$ to $x^2/d$. A fixed point would be
$d=x$, which would require $a\mid2x$, impossible because $(a,x)=1$
and $a\ge3$ is odd. The fibers therefore have size exactly two.
In primitive variables this is the familiar recovery
$h=(u+v)/a$ and $\{y,z\}=\{pDvh,pDuh\}$.
\end{proof}

These criteria are equivalent to established fixed-$x$ divisor
parametrizations; compare Bradford~\cite[Propositions~1--4]{Bradford}.
For his Type~I convention the divisor is $x^2/d$ in our notation.


\subsection{Mixing in the complete prime-power box}
\begin{lemma}[Full divisor-box mixing]\label{lem:box}
Let $a,x$ be positive integers with $(a,x)=1$, and let $m\ge0$.
Suppose that for every nonprincipal character
$\chi\bmod a$, at least $m$ distinct prime divisors $q\mid x$ satisfy
$\Re\chi(q)\le0$. Then, for every unit class $b\bmod a$,
\begin{equation}\label{eq:box}
 \left|\frac{\varphi(a)C_{a,x}(b)}{\tau(x^2)}-1\right|
 \le(\varphi(a)-1)3^{-m}.
\end{equation}
The counted prime divisors may be drawn from any chosen subwindow.
All additional primes and prime powers are included in $C_{a,x}$.
\end{lemma}
\begin{proof}
For $|z|=1$ and $e\ge1$ let
\[
 A_e(z)=\frac1{2e+1}\sum_{j=-e}^e z^j.
\]
Always $|A_e(z)|\le1$. If $\Re z\le0$, then for $e=1$,
$|A_1(z)|=|1+2\Re z|/3\le1/3$.
For $e\ge2$, write $z=e^{it}$. Since
$|\sin(t/2)|\ge1/\sqrt2$, the finite geometric sum gives
\[
 |A_e(z)|
 =\frac{|\sin((2e+1)t/2)|}{(2e+1)|\sin(t/2)|}
 \le\frac{\sqrt2}{2e+1}\le\frac{\sqrt2}{5}<\frac13.
\]
Using $x=\prod q^{e_q}$, character orthogonality and the ratio exponents
from Lemma~\ref{lem:dictionary} give
\[
 \frac{\varphi(a)C_{a,x}(b)}{\tau(x^2)}
 =\sum_{\chi\bmod a}\overline{\chi(b)}
              \prod_{q\mid x}A_{e_q}(\chi(q)).
\]
The principal term is $1$. Every nonprincipal term has absolute value
at most $3^{-m}$, by the hypothesis and the preceding factor bounds.
Summing proves \eqref{eq:box}.
\end{proof}

\begin{remark}
A prime power $q^e$ contributes $2e+1$ possible ratio exponents.
The normalized block $A_e$ accounts for all of them; replacing it by
$e$ independent ternary choices would instead count $3^e$ assignments.
\end{remark}

\section{A common good set for the growing catalogue}\label{sec:good}
\subsection{Parameters and sector tests}
For sufficiently large $N$, put
\begin{equation}\label{eq:parameters}
\begin{gathered}
 L=\log\log N,\qquad c=2^{-14},\qquad r=\lfloor cL\rfloor,\qquad k=2r,\\
 Q=\exp((\log N)^{1/4}),\qquad Z=N^{1/(8r)},\\
 \cP_N=\{N<p\le2N:p\text{ prime},\ p\equiv1\pmod4\},\\
 \cA_N=\{3\le a\le2(\log N)^\delta:a\equiv3\pmod4\}.
\end{gathered}
\end{equation}
For $p\in\cP_N$, $a\in\cA_N$, and nonprincipal $\chi\bmod a$, let
\[
 x_a=\frac{p+a}{4},\quad
 \cQ_a(p)=\{Q<q\le Z:q\text{ prime},\ q\mid x_a\},\quad
 K_a(p)=|\cQ_a(p)|,
\]
\[
 Z_{a,\chi}(p)=\#\{q\in\cQ_a(p):\operatorname{Re}\chi(q)\le0\}.
\]
Define the sufficient-test bad set
\begin{equation}\label{eq:badset}
 \cB_N=
 \bigcup_{a\in\cA_N}\{p\in\cP_N:K_a(p)<L/2\}
 \ \cup\!
 \bigcup_{\substack{a\in\cA_N\\\chi\ne\chi_0\ (a)}}
 \{p\in\cP_N:Z_{a,\chi}(p)<L/8\}.
\end{equation}
The complement of $\cB_N$ will be the common good set.

\begin{lemma}[Size of a character sector]\label{lem:sector}
Let $G$ be a finite abelian group and $\chi$ a nonprincipal character.
Then
\[
 S_\chi=\{g\in G:\Re\chi(g)\le0\},\qquad |S_\chi|\ge|G|/3.
\]
\end{lemma}
\begin{proof}
The image of $\chi$ consists of the $d$th roots of unity for some
$d\ge2$, each with equally many preimages. The number of roots in the
closed left half-plane is
\[
 \lfloor3d/4\rfloor-\lceil d/4\rceil+1.
\]
For $d\ge6$ this is at least $d/2-1\ge d/3$; for $d=2,3,4,5$ the
counts are $1,2,3,2$, respectively.
\end{proof}

\subsection{Distribution inputs and exact comparison means}
In the arithmetic sums in this section and their subsequent moment
application, the index $q$ always ranges over primes.
We record the precise standard distribution inputs. The
Siegel--Walfisz theorem implies, for every fixed $H>0$,
\begin{equation}\label{eq:sw}
 \pi(t;a,b)=\frac{\Li(t)}{\varphi(a)}
                +O_H\left(\frac{t}{(\log t)^H}\right)
 \quad ((a,b)=1,\ a\le\log t).
\end{equation}
Here a fixed Siegel--Walfisz modulus-range parameter $B=2$ suffices
in partial summation; its exponential error absorbs every fixed power
$(\log t)^{-H}$. The implied constant and onset may depend on $H$.
We also use the following consequence of the ordinary
Bombieri--Vinogradov theorem: for every fixed $A>0$,
\begin{equation}\label{eq:bv}
 \sum_{m\le4N^{1/4}}\max_{(v,m)=1}
 \left|\pi(2N;m,v)-\pi(N;m,v)
 -\frac{\Li(2N)-\Li(N)}{\varphi(m)}\right|
 \ll_A\frac{N}{(\log N)^A}.
\end{equation}
These forms follow from the character version of Siegel--Walfisz on
p.~2 and the two-term Bombieri--Vinogradov bound on p.~3 of
Vaughan's account \cite{VaughanBV}. Appendix~\ref{app:inputs} gives the
precise partial-summation steps and parameter substitutions.


For sufficiently large $N$, the parameters satisfy
\begin{equation}\label{eq:range}
 2(\log N)^\delta<Q<Z<N<p\quad(p\in\cP_N).
\end{equation}
Indeed, $Q<Z$ reduces to $8r<(\log N)^{3/4}$. Also $(a,x_a)=1$,
because a common divisor would divide $p$, while $a<p$.
All window primes are thus units modulo $a$ and exceed $a$.

For each unit class $b\pmod a$, put
\[
 \lambda_b=\sum_{\substack{Q<q\le Z\\q\equiv b\ (a)}}\frac1{q-1}.
\]
Partial summation in \eqref{eq:sw}, with $H=9$, gives uniformly in
$a\in\cA_N$ and $b$,
\begin{equation}\label{eq:lambda}
 \lambda_b=\frac{\log\log Z-\log\log Q}{\varphi(a)}
 +O((\log Q)^{-8})+O(Q^{-1}).
\end{equation}
To justify the uniformity, $a\le\log t$ holds eventually for all $t\ge Q$.
Replacing $1/q$ by $1/(q-1)$ costs $O(Q^{-1})$. The boundary and integral
errors in partial summation are bounded by a constant times
$(\log Q)^{-8}$. For a fixed $(a,\chi)$, the sector mean sums
$\lambda_b$ over $b\in S_\chi$; the total mean sums over all unit classes.
Thus their errors are at most
\[
 O\!\left(\varphi(a)((\log Q)^{-8}+Q^{-1})\right)
 =O\!\left((\log N)^{\delta-2}+(\log N)^\delta/Q\right)=o(1),
\]
uniformly in $a,\chi$. This is a sum within one statistic, distinct from
the later union over different moduli and characters.

For a nonprincipal character $\chi\pmod a$, define actual statistics
\begin{equation}\label{eq:statistics}
 I_q(p)=\one_{q\mid p+a},\qquad
 Z_{a,\chi}(p)=\sum_{\substack{Q<q\le Z\\q\bmod a\in S_\chi}}I_q(p),
 \qquad K_a(p)=\sum_{Q<q\le Z}I_q(p).
\end{equation}
The congruence $q\mid p+a$ is equivalent to $q\mid x_a$.
The exact comparison means are sums of $w_q=1/(q-1)$ over the respective
sets of primes, denoted $\mu_{a,\chi}$ and $\mu_{a,\mathrm{all}}$.
They are never replaced by asymptotic expressions inside a high moment.

The difference of logarithms in \eqref{eq:lambda} is
\begin{equation}\label{eq:log-window}
 \log\log Z-\log\log Q=\frac34L-\log(8r).
\end{equation}
The second term is not bounded as $N$ grows, but it is $o(L)$.
Writing $\sigma_\chi=|S_\chi|/\varphi(a)\ge1/3$, we obtain
\[
 \mu_{a,\chi}=\sigma_\chi\left(\frac34L-\log(8r)\right)+o(1),
 \qquad
 \mu_{a,\mathrm{all}}=\frac34L-\log(8r)+o(1).
\]
Consequently, for all sufficiently large $N$, uniformly in the parameters,
\begin{equation}\label{eq:mean-bounds}
 \frac{3L}{16}\le\mu_{a,\chi}\le L,
 \qquad \frac{5L}{8}\le\mu_{a,\mathrm{all}}\le L.
\end{equation}
Thus either $Z_{a,\chi}(p)<L/8$ or $K_a(p)<L/2$ forces deviation
from its own exact mean by at least $t_0=L/16$.

\subsection{A moment bound with explicit order dependence}
\begin{lemma}\label{lem:moment}
Let $B_i$ be finitely many independent Bernoulli variables of means
$w_i$, and put $Y=\sum_i(B_i-w_i)$ and $\mu=\sum_iw_i$.
For every integer $r\ge1$,
\begin{equation}\label{eq:moment}
 \E|Y|^{2r}\le2(4\mu)^r r!+2\,4^{2r}(2r)!
             \le2(4\mu r)^r+2(8r)^{2r}.
\end{equation}
\end{lemma}
\begin{proof}
For $|s|\le1$, independence and $\log(1+u)\le u$ yield
\[
 \log\E e^{sY}\le\mu(e^s-1-s)\le\mu s^2.
\]
When $\mu=0$ the claim is immediate. Otherwise use exponential Markov
on each tail, choosing $|s|=\min(t/(2\mu),1/2)$. This gives
\[
 \Prob(|Y|\ge t)\le2\exp\left\{-\min\left(\frac{t^2}{4\mu},
                                                   \frac t4\right)\right\}.
\]
Insert the bound into $2r\int_0^\infty t^{2r-1}\Prob(|Y|\ge t)\,dt$
and use $e^{-\min(u,v)}\le e^{-u}+e^{-v}$. The Gaussian and exponential
integrals give respectively $2(4\mu)^r r!$ and $2\,4^{2r}(2r)!$.
Finally $r!\le r^r$ and $(2r)!\le(2r)^{2r}$.
\end{proof}
For $\mu\le L$, $r\le cL$ and $t_0=L/16$, this implies
\begin{align}
 \frac{\E|Y|^{2r}}{t_0^{2r}}
 &\le2(1024r/L)^r+2(128r/L)^{2r}\notag\\
 &\le4\,2^{-4r}
 \le4(\log N)^{-2c\log2},\label{eq:normalized-moment}
\end{align}
where the last inequality uses $r\ge cL/2$, valid once $cL\ge2$.
Every constant in this estimate is independent of the growing order $r$.

\subsection{The coefficient of a distinct support}
\begin{lemma}\label{lem:compression}
Let $\mathcal Q$ be any finite set, $I_q\in\{0,1\}$,
$T=\sum_{q\in\mathcal Q}I_q$, and let $k$ be a nonnegative integer.
Zero powers are interpreted as $1$, including at zero, as polynomial
identities; thus the case $k=0$ is included.
The multilinear expansion of $(T-\mu)^k$ is
\begin{equation}\label{eq:multilinear}
 (T-\mu)^k=\sum_{\substack{S\subseteq\mathcal Q\\|S|\le k}}
                  \alpha_S\prod_{q\in S}I_q,
\end{equation}
where, for $s=|S|$,
\begin{equation}\label{eq:alpha}
 \alpha_S=\sum_{j=s}^k\binom{k}{j}(-\mu)^{k-j}s!
                      \left\{\begin{matrix}j\\s\end{matrix}\right\}.
\end{equation}
For $\mu\ge0$, $|\alpha_S|\le(\mu+s)^k$. The empty-support
coefficient is $(-\mu)^k$.
\end{lemma}
\begin{proof}
Expand the power and use $I_q^j=I_q$ for $j\ge1$. In a product of
$j$ indicators, using exactly the prescribed $s$ indices is counted by
the number of surjections onto those indices, namely
$s!\left\{\begin{smallmatrix}j\\s\end{smallmatrix}\right\}$.
This proves \eqref{eq:alpha}. That number is at most $s^j$; applying the
binomial theorem to the absolute values gives the asserted bound.
For $s=0$, only the term $j=0$ remains, with the usual convention
$\left\{\begin{smallmatrix}0\\0\end{smallmatrix}\right\}=1$.
\end{proof}
For our parameters $k=2r$, $s\le2r$, and $\mu\le L$, normalization
before estimating gives the uniform bound
\begin{equation}\label{eq:coefficient-budget}
 \frac{|\alpha_S|}{t_0^{2r}}
 \le[16(1+2c)]^{2r}\le32^{2r}
 \le(\log N)^{\beta},\qquad \beta=10c\log2<1.
\end{equation}
It is this bound per distinct support, rather than a bound on the
number of ordered tuples, that permits use of ordinary
Bombieri--Vinogradov with $r$ growing.

\subsection{Transfer by distinct supports}
Fix $a\in\cA_N$ and a set $\mathcal Q$ of primes with
$a<q\le Z$, all odd. Write $T(p)=\sum_{q\in\mathcal Q}I_q(p)$
and assume that its exact comparison mean is at most $L$.
This includes the statistics in \eqref{eq:statistics}.
In the comparison model,
let $B_q$ be independent with mean $1/(q-1)$, let $T^*=\sum B_q$,
and let $\mu$ be the exact sum of their means. For each support $S$
in Lemma~\ref{lem:compression}, put $d(S)=\prod_{q\in S}q$.
The simultaneous divisibility event for actual $p\in\cP_N$ is exactly
\[
 p\equiv1\pmod4,\qquad p\equiv-a\pmod{d(S)}.
\]
Since the primes $q$ exceed $a$, this is one reduced class modulo
$4d(S)$. Write $P_0=(\Li(2N)-\Li(N))/2$. Then
\begin{equation}\label{eq:joint-count}
 \sum_{p\in\cP_N}\prod_{q\in S}I_q(p)
       =\frac{P_0}{\varphi(d(S))}+R_a(d(S)).
\end{equation}
The model joint expectation is exactly $1/\varphi(d(S))$.
For the empty support, $d(S)=1$ and the remainder is
$|\cP_N|-P_0$; it is included in this formula.

Unique prime factorization ensures that distinct supports give distinct
$d(S)$. Every such support has at most $2r$ primes, so
\begin{equation}\label{eq:modulus-budget}
 4d(S)\le4Z^{2r}=4N^{1/4}.
\end{equation}
Using \eqref{eq:bv} with $A=5$, the sum of absolute remainders over
these supports is $O(N/(\log N)^5)$. In particular, it is not necessary
to multiply this bound by the number of supports.
Equations \eqref{eq:multilinear}--\eqref{eq:joint-count} and
\eqref{eq:coefficient-budget} yield
\begin{equation}\label{eq:transfer}
 \frac1{t_0^{2r}}\sum_{p\in\cP_N}(T(p)-\mu)^{2r}
 =P_0\frac{\E(T^*-\mu)^{2r}}{t_0^{2r}}
       +O\bigl(N(\log N)^{-5+\beta}\bigr).
\end{equation}
The signed expansion is used only to establish the equality and its
error bound; the actual even moment on the left is nonnegative.

By Markov's inequality and \eqref{eq:normalized-moment},
any event forcing $|T(p)-\mu|\ge t_0$ occurs for at most
\begin{equation}\label{eq:one-bad}
 4P_0\,2^{-4r}+O\bigl(N(\log N)^{-5+\beta}\bigr)
\end{equation}
primes. In particular, this applies to each failed threshold in
\eqref{eq:badset}, by \eqref{eq:mean-bounds}.

\begin{proposition}[Common exceptional-set estimate]\label{prop:commonbad}
For the parameters in \eqref{eq:parameters},
\[
 |\cB_N|\ll|\cP_N|(\log N)^{-\delta}.
\]
\end{proposition}
\begin{proof}
For each $a$ we have $\varphi(a)-1$ sector statistics and one total
factor statistic. The sufficient-test bad set is $\cB_N$ from \eqref{eq:badset}.
Its constituent events need not be disjoint; we only use a union bound.
The number of statistics over all $a\in\cA_N$ is at most
\begin{equation}\label{eq:union-budget}
 \sum_{a\le2(\log N)^\delta}\varphi(a)\ll(\log N)^{2\delta}.
\end{equation}
A union bound in \eqref{eq:one-bad}, using
$P_0\asymp|\cP_N|\asymp N/\log N$, shows that the proportion of
bad primes is
\begin{equation}\label{eq:bad-rate}
 O\left((\log N)^{2\delta-2c\log2}
        +(\log N)^{-4+\beta+2\delta}\right)
       =O((\log N)^{-\delta}).
\end{equation}
For the equality, note $c=64\delta$, so
$2c\log2>3\delta$ and $4-\beta>3\delta$.

\end{proof}

\section{Proofs of the main results}
\begin{proof}[Proof of Theorem~\ref{thm:main}]
For $p\in\cP_N\setminus\cB_N$, every nonprincipal character modulo
every $a\in\cA_N$ has at least $L/8$ distinct prime divisors of $x_a$
in its sector. Lemma~\ref{lem:box} therefore gives, for every unit class,
\begin{equation}\label{eq:full-relative}
 \left|\frac{\varphi(a)C_{a,x_a}(b)}{\tau(x_a^2)}-1\right|
 \le(\varphi(a)-1)3^{-L/8}
 \le2(\log N)^{\delta-\log3/8}
 =2(\log N)^{-\kappa}.
\end{equation}
All additional prime factors and their powers are already included in
$C_{a,x_a}$; no information about them beyond Lemma~\ref{lem:box} is
needed. Proposition~\ref{prop:exact}, at $b=-p^{-1}$ and $b=-1$,
now gives the two complete local counts. The first target depends on
$p$, which is harmless because \eqref{eq:full-relative} holds for all
unit classes for each individual good prime.

For sufficiently large $N$ and $N<p\le2N$,
$(\log p)^\delta\le2(\log N)^\delta$. The prescribed shifts
$a_j=3+4j$, $0\le j<J(p)$, are exactly those satisfying
$3\le a_j\le(\log p)^\delta$, so they belong to $\cA_N$.
Also $\log p\asymp\log N$, and hence the relative errors are
$O((\log p)^{-\kappa})$, uniformly in $j$.
They are eventually less than $1/2$, so the counts in the ratio are
positive. Taking their ratio and their sum gives
\eqref{eq:ratio-total}. Strict minimality distinguishes solutions at
different positions.

Choose a fixed sufficiently large dyadic starting point and let
$\mathcal E$ be the union of the sets $\cB_N$ for disjoint dyadic
intervals above it, together with the finite initial range of primes.
This is a fixed sufficient exceptional set. To estimate its counting
function up to $X$, discard primes at most $\sqrt X$; their number is
negligible compared with $X/(\log X)^{1+\delta}$. For the remaining
dyadic intervals $N\gg\sqrt X$, the logs are uniformly comparable to
$\log X$. Proposition~\ref{prop:commonbad} and the prime number theorem give
\[
 \sum_N|\cB_N|\ll(\log X)^{-\delta}
                      \sum_N|\cP_N|
 \ll X/(\log X)^{1+\delta}.
\]
A possible final partial interval is bounded by its full dyadic
interval. This proves \eqref{eq:global-exception}.

\end{proof}

\subsection{A broader total-factor window}
Keep the sector window unchanged, and also put
\begin{equation}\label{eq:wide-definition}
 W=2(\log N)^\delta,\qquad
 \widetilde K_a(p)=\sum_{W<q\le Z}I_q(p),\qquad
 \widetilde\mu=\sum_{W<q\le Z}\frac1{q-1}.
\end{equation}
Every prime in this wider window exceeds every $a\in\cA_N$.
Its exact comparison mean does not depend on $a$.
The case $a=1$ of \eqref{eq:sw}, followed by partial summation, gives
\[
 \sum_{W<q\le Z}\frac1{q-1}
 =\int_W^Z\frac{dt}{t\log t}
       +O((\log W)^{-2}+W^{-1}).
\]
For example, $H=3$ suffices for the integrated prime-counting error;
replacing $1/q$ by $1/(q-1)$ costs $O(W^{-1})$.
Since $W\to\infty$, this proves
\begin{equation}\label{eq:wide-mean}
 \widetilde\mu
 =L-\log(8r)-\log(\log2+\delta L)+o(1)
 =L-O(\log L).
\end{equation}
In particular, $\widetilde\mu\le L$ eventually. For every fixed
$0<\theta<15/16$ it is eventually at least $(\theta+1/16)L$.
Thus $\widetilde K_a(p)<\theta L$ forces a deviation of at least
$t_0=L/16$ from its own exact mean. The lower endpoint requires only the ordinary prime number theorem.

\begin{proposition}[An enlarged sufficient-test bad set]\label{prop:wide}
Fix $0<\theta<15/16$ and define
\begin{equation}\label{eq:wide-badset}
 \cB_N^*=\cB_N\cup
 \bigcup_{a\in\cA_N}
       \{p\in\cP_N:\widetilde K_a(p)<\theta L\}.
\end{equation}
Then, for sufficiently large $N$ depending on $\theta$,
\[
 |\cB_N^*|\ll_\theta |\cP_N|(\log N)^{-\delta}.
\]
\end{proposition}
\begin{proof}
Equation~\eqref{eq:wide-mean} gives $\widetilde\mu\le L$ and the
required deviation $L/16$. Thus the normalized model moment is at most
$4\,2^{-4r}$ by \eqref{eq:normalized-moment}. Every distinct support
still has normalized coefficient at most $(\log N)^\beta$ by
\eqref{eq:coefficient-budget}, and its modulus is at most
$4Z^{2r}=4N^{1/4}$ by \eqref{eq:modulus-budget}. The lower endpoint has
changed, but all support primes exceed $a$ and are odd. Consequently
the same reduced-class count and the same unweighted BV remainder sum
apply. Formula~\eqref{eq:one-bad} bounds the number of added failures
for each $a$.

There are $|\cA_N|\ll(\log N)^\delta$ new statistics, compared with
the $O((\log N)^{2\delta})$ original statistics. The enlarged union
therefore still has bad proportion
\[
 O\!\left((\log N)^{2\delta-2c\log2}
          +(\log N)^{-4+\beta+2\delta}\right)
 =O((\log N)^{-\delta}).
\]
Here $c=64\delta$, $2c\log2>3\delta$, and $4-\beta>3\delta$.
The necessary onset, and the constant after absorbing an initial
range, may depend on the fixed $\theta$.
\end{proof}

\begin{proof}[Proof of Corollary~\ref{cor:many}]
On $\cP_N\setminus\cB_N^*$ there are at least $\theta L$ distinct
prime factors of $x_a$ in the wider window. Each contributes at least
$3$ to $\tau(x_a^2)$, so
\[
 \tau(x_a^2)\ge3^{\widetilde K_a(p)}
                 \ge(\log N)^{\theta\log3}.
\]
The relative error in \eqref{eq:full-relative} is eventually at most
$1/2$. The exact counts then give dyadic lower bounds with denominators
$2\varphi(a)$ for Type~I and $4\varphi(a)$ for Type~II.
For fixed $\theta$, converting $\log N$ to $\log p$ costs at most
another factor two for sufficiently large $N$. This proves
\eqref{eq:many} on the enlarged good set. Define $\mathcal E_\theta$
from the dyadic $\cB_N^*$ and a sufficiently large finite initial range,
including $\mathcal E$. The same dyadic summation as above proves the
claimed counting bound. The original asymptotics remain valid because
the enlarged bad set contains the original one.
\end{proof}

\begin{proof}[Proof of Corollary~\ref{cor:fixed}]
For fixed $J_0$, the prescribed positions are included once
\[
 p\ge P_{J_0}:=\exp((4J_0-1)^{1/\delta}),
\]
since this is exactly the condition $J(p)\ge J_0$.
Theorem~\ref{thm:main}, respectively Corollary~\ref{cor:many}, applies
above this threshold. The finitely many smaller primes are absorbed
in the constants, which are allowed to depend on $J_0$ and, for the
multiplicity assertion, on $\theta$. No effective starting point is
obtained by this deduction.
\end{proof}

\section{Arithmetic interpretation and parameters}
\subsection{Quadratic reciprocity}
The following observation interprets one of the sector tests already
used. Here $p\equiv1\pmod4$ is prime, $a\equiv3\pmod4$,
$3\le a<p$, $x=(p+a)/4$, and $x=\prod q^{e_q}$,
$d=\prod q^{j_q}$ in the notation of Lemma~\ref{lem:dictionary}.

\begin{proposition}[Reciprocity interpretation]\label{prop:resource}
For every prime factor $q\mid x$,
\begin{equation}\label{eq:reciprocity}
 \left(\frac q p\right)=\left(\frac q a\right),
\end{equation}
where the right hand side is the Jacobi symbol. The character
$\chi_a(n)=(n/a)$ on the unit group modulo $a$ is nonprincipal and odd.
Thus the corresponding sector count is exactly
\[
 Z_{a,\chi_a}(p)=
 \#\{Q<q\le Z:q\mid x,\ q\text{ prime},\ (q/p)=-1\}.
\]
For each Type~II witness $d\mid x^2$, the primitive resource is
\[
 R=uv=\frac{dx}{(d,x)^2}
      =\prod_{q\mid x}q^{|j_q-e_q|}\mid x.
\]
It satisfies $(R/p)=-1$, and in particular
$R\ge n_-(p)$, the least positive quadratic nonresidue modulo $p$.
\end{proposition}
\begin{proof}
For odd $q\mid x$, $p\equiv-a\pmod q$ and quadratic reciprocity give
\[
 (q/p)=(p/q)=(-a/q)
 =(-1/q)(a/q)=(q/a),
\]
because $p\equiv1\pmod4$ and $a\equiv3\pmod4$ make the two signs
in the last step cancel. If $q=2$, then $p\equiv-a\pmod8$, so the
supplementary law gives $(2/p)=(2/a)$.
The value $\chi_a(-1)=-1$ proves that $\chi_a$ is odd and nonprincipal.
Its values on primes dividing $x$ are $\pm1$, so its sector is exactly
the negative-value set just displayed.

The formula for $R$ follows prime by prime from
Lemma~\ref{lem:dictionary}. In Type~II, $u/v\equiv-1\pmod a$.
As Jacobi values on units are $\pm1$,
\[
 \chi_a(uv)=\chi_a(u)\chi_a(v)^{-1}
          =\chi_a(uv^{-1})=\chi_a(-1)=-1.
\]
Multiplicativity and \eqref{eq:reciprocity} now give $(R/p)=-1$.
The least-nonresidue bound follows from its definition.
\end{proof}

This identifies the quadratic-character sector condition with
prime-specific arithmetic: on the sufficient good set, every prescribed
$x_a$ has at least $L/8$ distinct nonresidue prime factors in the
specified window. The resource $uv$ is a divisor of $x_a$ attached to the primitive
ratio of a Type~II witness.

\subsection{Parameter budget}
The catalogue and exceptional-set exponents can be separated. Keep
$c=2^{-14}$, $r=\lfloor cL\rfloor$, and $\beta=10c\log2$.
For fixed $\eta,\sigma>0$, a sufficient budget for a catalogue
$a\le2(\log N)^\eta$ and an exceptional-proportion saving
$(\log N)^{-\sigma}$ is
\begin{equation}\label{eq:separate-budget}
\begin{gathered}
 \eta<\min\{1/4,\ (\log2)/16,\ (\log3)/8\},\\
 2\eta+\sigma<2c\log2,\qquad 2\eta+\sigma<4-\beta.
\end{gathered}
\end{equation}
The first line makes the errors in the uniform mean estimates and both
auxiliary Fourier estimates tend to zero. The full-box estimate then
has relative error $O((\log N)^{\eta-\log3/8})$.
There are $O((\log N)^{2\eta})$ original statistics; the second line
pays for their union in the moment and BV estimates. The broader
window $W=2(\log N)^\eta$ adds only $O((\log N)^\eta)$ statistics,
and its mean is still $L-O_\eta(\log L)$. Its threshold argument works
for every fixed $\theta<15/16$. Constants and onsets may depend on
these fixed parameters. The present choice $\eta=\sigma=\delta$
satisfies these inequalities with room to spare. This is a sufficient
budget, not an optimization.

The all-class estimate supplies both targets $-1$ and $-p^{-1}$
without another union over residue classes. The resulting $2:1$ ratio
holds on the common good set. It can fail at an individual position:
$(p,a,x)=(53,7,15)$ has $\FI=2$ and $\FII=0$.

The supplementary exact checks cover the divisor dictionary, complete
local counts, prime-power boxes, recovery multiplicities, and the
coefficient budgets. Their finite ranges and arithmetic conventions
are recorded with the verification code.

\appendix
\section{Squarefree subfamilies and one-sided recovery}\label{app:auxiliary}
The complete-count argument also supplies useful squarefree
subfamilies. Three-way allocation of coprime factors, and its $3^K$
count, already occurs in Browning and Elsholtz \cite[Section~2]{BE}
and Elsholtz and Planitzer \cite[Section~7]{EP}.

\subsection{Binary and ternary character estimates}
\begin{lemma}[Binary subset products]\label{lem:binary}
Let $G$ be a finite abelian group, $K\in\mathbb Z_{\ge0}$, $m\ge0$,
and $g_1,\ldots,g_K\in G$. Suppose that every nonprincipal character
has at least $m$ indices $i$ with $\Re\chi(g_i)\le0$.
If $D(b)$ counts subsets with product $b$, then
\begin{equation}\label{eq:fourier-bound}
 \left|\frac{|G|D(b)}{2^K}-1\right|
 \le(|G|-1)2^{-m/2}\qquad(b\in G).
\end{equation}
If the right-hand side is at most $1/2$, then
$D(b)\ge2^{K-1}/|G|$ for every $b\in G$.
\end{lemma}
\begin{proof}
Let $\nu(b)=2^{-K}D(b)$, and use the Fourier convention
\[
 \widehat\nu(\chi)=\sum_{b\in G}\nu(b)\chi(b),\qquad
 \nu(b)=\frac1{|G|}\sum_{\chi\in\widehat G}
                \widehat\nu(\chi)\overline{\chi(b)}.
\]
Expanding over the choice of each index gives the exact identity
\[
 \widehat\nu(\chi)=\prod_{i=1}^K\frac{1+\chi(g_i)}2.
\]
Every factor has absolute value at most $1$. On the specified $m$ indices,
\[
 \left|\frac{1+\chi(g_i)}2\right|^2
     =\frac{1+\operatorname{Re}\chi(g_i)}2\le\frac12.
\]
The nonprincipal Fourier coefficients therefore have absolute value at
most $2^{-m/2}$. The principal coefficient is exactly $1$, so Fourier
inversion and the triangle inequality give \eqref{eq:fourier-bound}.
For the trivial group the nonprincipal sum is empty and the error is zero.
When $K=0$, the only subset is the empty subset; the same formulas apply.
\end{proof}

\begin{lemma}\label{lem:ternary}
Let $G$ be a finite abelian group, $K\in\mathbb Z_{\ge0}$, $m\ge0$, and
$g_1,\ldots,g_K\in G$. For $b\in G$, let $T(b)$ count assignments
$\epsilon_i\in\{-1,0,1\}$ with $\prod_i g_i^{\epsilon_i}=b$.
If every nonprincipal character has at least $m$ indices $i$ with
$\operatorname{Re}\chi(g_i)\le0$, then
\begin{equation}\label{eq:ternary-fourier-bound}
 \left|\frac{|G|T(b)}{3^K}-1\right|
 \le (|G|-1)3^{-m}.
\end{equation}
In particular, if the right-hand side is at most $1/2$, then
$T(b)\ge3^K/(2|G|)$ for every $b\in G$.
\end{lemma}
\begin{proof}
The normalized character transform of the three choices at index $i$ is
\[
 \frac{1+\chi(g_i)+\chi(g_i)^{-1}}3
 =\frac{1+2\operatorname{Re}\chi(g_i)}3.
\]
It has absolute value at most $1$ in general and at most $1/3$ when
$\operatorname{Re}\chi(g_i)\le0$. Hence every nonprincipal coefficient has absolute value
at most $3^{-m}$, while the principal coefficient is $1$. Fourier inversion
proves \eqref{eq:ternary-fourier-bound}; the trivial-group and $K=0$
cases are immediate.
\end{proof}

\subsection{Simultaneous window supply}
For a unit class $b\pmod a$, let $D_{a,p}(b)$ count subsets of
$\cQ_a(p)$ whose product is $b$ modulo $a$, with the empty product equal
to $1$. Let $T_{a,p}(b)$ count ordered pairs $(u,v)$ obtained as products
of disjoint subsets of $\cQ_a(p)$ with $u/v\equiv b\pmod a$.
The ratio is taken in the unit group. There are three choices per
window prime, and every such pair satisfies $(u,v)=1$ and $uv\mid x_a$.

\begin{theorem}[Auxiliary squarefree-window supply]\label{thm:supply}
With \eqref{eq:parameters} and \eqref{eq:badset},
\[
 |\cB_N|\ll|\cP_N|(\log N)^{-\delta}.
\]
Every $p\in\cP_N\setminus\cB_N$ satisfies, simultaneously for all
$a\in\cA_N$ and all unit classes $b\pmod a$,
\begin{equation}\label{eq:supply}
\begin{gathered}
 K_a(p)\ge L/2,\qquad
 D_{a,p}(b)\ge\frac{2^{K_a(p)-1}}{\varphi(a)},\\
 T_{a,p}(b)\ge\frac{3^{K_a(p)}}{2\varphi(a)}
 \ge\frac{(\log N)^{\gamma_0}}{2\varphi(a)},
 \qquad \gamma_0=\frac{\log3}{2}.
\end{gathered}
\end{equation}
\end{theorem}
\begin{proof}
The exceptional-set estimate is Proposition~\ref{prop:commonbad}.
For every remaining $p$, all $a\in\cA_N$ simultaneously satisfy
\[
 K_a(p)\ge L/2,\qquad
 Z_{a,\chi}(p)\ge L/8\quad(\chi\ne\chi_0).
\]
Apply Lemma~\ref{lem:binary} and Lemma~\ref{lem:ternary} to the residues
of the distinct primes in $\cQ_a(p)$. The binary relative error is at most
\[
 (\varphi(a)-1)2^{-L/16}
 \le2(\log N)^{\delta-\log2/16}=o(1),
\]
while the ternary relative error is at most
\[
 (\varphi(a)-1)3^{-L/8}
 \le2(\log N)^{\delta-\log3/8}=o(1).
\]
Both are eventually at most $1/2$, uniformly in $a$. Thus every unit class
satisfies both the displayed binary bound and
$T_{a,p}(b)\ge3^{K_a(p)}/(2\varphi(a))$. Since
$3^{L/2}=(\log N)^{\gamma_0}$, the last ternary lower bound in
\eqref{eq:supply} follows. The selected primes are distinct actual divisors
of $x_a$; this verifies the integer-divisor interpretation as well as both
residue-class counts.
\end{proof}

The exponent $\gamma_0=(\log3)/2$ concerns this selected window;
Corollary~\ref{cor:many} uses the broader total-factor window.

\subsection{A one-sided Type~II construction}
\begin{lemma}\label{lem:recovery}
Let $p$ be prime, $a\equiv3\pmod4$, $3\le a<p$, and $p\equiv1\pmod4$.
Put $x=(p+a)/4$. If $d$ is a positive divisor of $x$ and
$d\equiv-1\pmod a$, then, with
\begin{equation}\label{eq:recovery}
 h=\frac{d+1}{a},\qquad y=\frac{phx}{d},\qquad z=phx,
\end{equation}
we have a Type~II solution $x<y<z$ of \eqref{eq:es}. Distinct such $d$
give distinct unordered triples at the fixed position $x$.
\end{lemma}
\begin{proof}
All quantities in \eqref{eq:recovery} are positive integers. Since
$d\equiv-1\pmod a$ and $a\ge3$, we have $d\ge a-1\ge2$. Moreover
\[
 \frac1x+\frac1y+\frac1z
 =\frac1x+\frac{d+1}{phx}
 =\frac{p+a}{px}=\frac4p.
\]
Both $y$ and $z$ are multiples of $p$ and at least $p$, whereas
$x<p/2$. Hence $x$ is strictly smallest, $p\nmid x$, and $z/y=d>1$.
The ratio of the larger to the smaller of the remaining two denominators
recovers $d$ from the unordered triple, proving injectivity.
\end{proof}

\section{Fixed catalogues from earlier inputs}\label{app:fixed}
The following consequence explains the distinction between the
fixed-catalogue precedent and the uniform quantitative application.

\begin{proposition}[Fixed-catalogue consequence]\label{prop:fixed-known}
Let $\mathcal A\subset\{3,7,11,\ldots\}$ be fixed and finite. On a relative-density-one set of primes
$p\equiv1\pmod4$, simultaneously for $a\in\mathcal A$ and
$x_a=(p+a)/4$,
\[
 \FI(p,x_a)\sim\frac{\tau(x_a^2)}{\varphi(a)},\qquad
 \FII(p,x_a)\sim\frac{\tau(x_a^2)}{2\varphi(a)}.
\]
For every fixed $0<\eps<\log3$, each count is also
$\gg_{\mathcal A,\eps}(\log p)^{\log3-\eps}$ on a possibly
smaller relative-density-one set.
\end{proposition}
\begin{proof}
For fixed $a$ and a nonprincipal character $\chi\bmod a$, define
the nonnegative strongly additive function
\[
 w_{a,\chi}(n)=\sum_{\substack{q\mid n\\q\nmid a}}
                      \one_{\Re\chi(q)\le0}.
\]
Let $\sigma_\chi=|S_\chi|/\varphi(a)$. Lemma~\ref{lem:sector}
gives $\sigma_\chi\ge1/3$. The prime number theorem in fixed
arithmetic progressions gives
\[
 \sum_{q\le N}\frac{w_{a,\chi}(q)}{q-1}
       =\sigma_\chi\log\log N+O_a(1),\qquad
 \max_{q\le N}w_{a,\chi}(q)\le1.
\]
The fixed-shift normal-order theorem of Barban and Elliott, recalled
as Lemma~A.3 of \cite{ET}, therefore applies to $p+a$.
After a finite union over $(a,\chi)$, its conclusion holds for all
$w_{a,\chi}(p+a)$ outside $o(\pi(X))$ primes up to $X$.
Since $\pi(X;4,1)\sim\pi(X)/2$, restricting to $p\equiv1\pmod4$
preserves relative density one. On this progression $x_a$ is an
integer, and dividing $p+a$ by $4$ removes at most one distinct prime
factor. Consequently, on a relative-density-one set in this progression,
\[
 w_{a,\chi}(x_a)\ge\tfrac14\log\log p
 \qquad(a\in\mathcal A,\ \chi\ne\chi_0).
\]
Lemma~\ref{lem:box} now yields, uniformly over
this fixed catalogue and all unit classes $b$,
\[
 C_{a,x_a}(b)=\frac{\tau(x_a^2)}{\varphi(a)}
 \left(1+O_{\mathcal A}((\log p)^{-(\log3)/4})\right).
\]
Proposition~\ref{prop:exact} gives the asserted complete counts.

Apply the same fixed-shift normal-order theorem to $w=\omega$.
Restricting to the same progression and dividing by $4$ as above gives,
for all $a\in\mathcal A$ simultaneously,
$\omega(x_a)=(1+o(1))\log\log p$ on a relative-density-one set there.
Since $\tau(x_a^2)\ge3^{\omega(x_a)}$, the claimed lower bounds
follow, with any fixed positive $\eps$ providing the exponent margin.
\end{proof}

This deduction uses finitely many fixed shifts and supplies an
$o(\pi(X))$ exceptional set. The common logarithmic-power estimate
for a specified growing catalogue requires the uniform calculation
in the main text or Appendix~\ref{app:elliott}.

\subsection{The per-class divisor route}
For fixed $J_0$ and $0<\eps<\log2$, the per-class divisor estimates in
the final part of the proof of Bettin et al.\
\cite[Lemma~4.9 and equation~(4.3)]{Bettin} also imply
\begin{equation}\label{eq:bettin-fixed}
 \FII\bigl(p,(p+3)/4+j\bigr)
 \gg_{J_0,\eps}(\log p)^{\log2-\eps}
 \qquad(0\le j<J_0)
\end{equation}
for a density-one set of the relevant primes. This is an extraction
from their proof, not an Erd\H{o}s--Straus theorem named there.
Indeed, for a fixed $a=3+4j$, take their shift $p_0=-a$, modulus
$4a$, and unit class $2a-1$. Their estimate supplies
$\gg_{a,\eps}(\log p)^{\log2-\eps}$ divisors $d\mid p+a$ in that class.
Such $d$ are odd, hence divide $x=(p+a)/4$, and satisfy $d\equiv-1\pmod a$.
Lemma~\ref{lem:recovery} converts different $d$ into different solutions.
A finite union over the fixed shifts proves \eqref{eq:bettin-fixed}.
This argument does not extract uniform constants for growing $J_0$ or
the explicit exceptional-set estimate \eqref{eq:global-exception}.

\section{An alternative transfer by Elliott's model}\label{app:elliott}
Elliott's Bernoulli approximation, in the explicit form of
\cite[Theorem~1.1]{Gorodetsky}, can replace the moment transfer.
We verify its hypotheses uniformly for the growing shifts.
In the form we use, let $V$ take values in $[1,T]\cap\mathbb N$,
let $2\le y\le T$, let $\mathcal Q$ be a set of primes at most $y$, and let
$g:\mathbb N\to[0,1]$ be multiplicative, vanish off
$\langle\mathcal Q\rangle$, and satisfy $g(q)<1$ for $q\in\mathcal Q$.
Let $B_q$ be independent Bernoulli variables with means $g(q)$ and put
$S=\sum_{q\le y}g(q)\log q/(1-g(q))$.
For $S>0$ and $\log z\ge8\max\{\log y,S\}$, Elliott's bound is
\begin{align}
 &d_{\mathrm{TV}}\bigl((\one_{q\mid V})_{q\in\mathcal Q},
                         (B_q)_{q\in\mathcal Q}\bigr)\notag\\
 &\quad\le10\exp\left(-\frac{\log z}{8\log y}
                              \log\frac{\log z}{S}\right)
 +12\!\sum_{\substack{d\le z^4\\d\in\langle\mathcal Q\rangle}}
       \mu^2(d)4^{\omega(d)}\bigl|\Prob(d\mid V)-g(d)\bigr|.
       \label{eq:elliott-input}
\end{align}
Keep \eqref{eq:parameters} and \eqref{eq:wide-definition}, and put
\[
\begin{gathered}
 y=Z,\qquad z=N^{1/16},\qquad H=|\cP_N|,\\
 P_0=\frac{\Li(2N)-\Li(N)}2,\qquad
 \mathcal Q_N=\{q\text{ prime}:W<q\le Z\}.
\end{gathered}
\]
For $a\in\cA_N$, choose $p$ uniformly from $\cP_N$; then
$x_a=(p+a)/4$ is an integer in $[1,N]$ for sufficiently large $N$.
Let $(B_q)_{q\in\mathcal Q_N}$ be independent Bernoulli variables
with $\Prob(B_q=1)=1/(q-1)$.

\begin{proposition}[Uniform Bernoulli approximation]\label{prop:elliott-uniform}
For every fixed $B>0$,
\begin{equation}\label{eq:elliott-tv}
 \sup_{a\in\mathcal A_N}
 d_{\mathrm{TV}}\!\left(
   (\one_{q\mid x_a})_{q\in\mathcal Q_N},
   (B_q)_{q\in\mathcal Q_N}\right)
 \ll_B(\log N)^{-B}.
\end{equation}
The implied constant is independent of the growing shift $a$.
\end{proposition}
\begin{proof}
Define the multiplicative function $g$ by $g(1)=1$,
$g(q^j)=1/\varphi(q^j)$ for $q\in\mathcal Q_N$, $j\ge1$, and
$g(d)=0$ outside the semigroup $\langle\mathcal Q_N\rangle$.
This takes values in $[0,1]$ and has $g(q)<1$ at every window prime.
The parameter in Elliott's theorem satisfies
\[
 S=\sum_{W<q\le y}\frac{\log q}{q-2}
 \le\log y+O(1)\le2\log y.
\]
Here Mertens' estimate for $\sum_{q\le y}(\log q)/q$ suffices;
the correction from $q$ to $q-2$ has a bounded sum over $q\ge3$.
Since $\log z/\log y=r/2$, the hypothesis
$\log z\ge8\max\{\log y,S\}$ holds once $r\ge32$.
The first term of Elliott's bound is then at most
\begin{equation}\label{eq:elliott-exponential}
 10\exp\!\left(-\frac r{16}\log\frac r4\right)
 \ll_B(\log N)^{-B}.
\end{equation}
The last assertion follows from $r\sim c\log\log N$.

It remains to bound the theorem's weighted remainder sum, whose range
is $d\le z^4=D=N^{1/4}$ and whose weight is
$\mu^2(d)4^{\omega(d)}$. Every $d\in\langle\mathcal Q_N\rangle$
is odd and coprime to every $a\in\mathcal A_N$. The congruences
$p\equiv1\pmod4$ and $p\equiv-a\pmod d$ specify one reduced class
modulo $4d$. Consequently, on writing
\[
 \#\{p\in\mathcal P_N:d\mid x_a\}
   =\frac{P_0}{\varphi(d)}+R_a(d),
\]
we have the exact normalization
\begin{equation}\label{eq:elliott-normalization}
 \Prob(d\mid x_a)-g(d)
 =\frac{R_a(d)}H+\frac{P_0-H}{H\varphi(d)}.
\end{equation}
Define
\[
 E_d=\max_{(b,4d)=1}
 \left|\pi(2N;4d,b)-\pi(N;4d,b)
       -\frac{\Li(2N)-\Li(N)}{\varphi(4d)}\right|.
\]
Then $|R_a(d)|\le E_d$ uniformly in $a$. Equation~\eqref{eq:bv} gives, for every fixed $A>0$
\begin{equation}\label{eq:elliott-bv}
 \sum_{d\le D}E_d\ll_A N(\log N)^{-A}.
\end{equation}
Indeed, the moduli $4d\le4N^{1/4}$ are strictly within its range.
The trivial count of integers in a residue class also gives
$E_d\ll N/\varphi(d)$ for $d\le D$.
For each fixed positive integer $j$, an Euler product and Mertens'
estimate yield
\begin{equation}\label{eq:elliott-euler}
 \sum_{d\le D}\frac{\mu^2(d)j^{\omega(d)}}{\varphi(d)}
 \le\prod_{q\le D}\left(1+\frac j{q-1}\right)
 \ll_j(\log(2D))^j.
\end{equation}
Cauchy--Schwarz, followed by \eqref{eq:elliott-bv} and
\eqref{eq:elliott-euler} with $j=16$, therefore gives
\begin{align}
 \sum_{\substack{d\le D\\d\in\langle\mathcal Q_N\rangle}}
       \mu^2(d)4^{\omega(d)}E_d
 &\le
 \left(\sum_{d\le D}\mu^2(d)16^{\omega(d)}E_d\right)^{1/2}
 \left(\sum_{d\le D}E_d\right)^{1/2}\notag\\
 &\ll_A N(\log N)^{8-A/2}.\label{eq:elliott-weighted-bv}
\end{align}

Since $H\asymp N/\log N$, its normalized contribution is
$O_A((\log N)^{9-A/2})$.
Also $|P_0-H|\le E_1\ll_A N(\log N)^{-A}$. Its contribution in
\eqref{eq:elliott-normalization}, using
\eqref{eq:elliott-euler} with $j=4$, is
$O_A((\log N)^{5-A})$.
Taking, for example, $A=2B+20$ bounds both remainder contributions
by $O_B((\log N)^{-B})$. Together with
\eqref{eq:elliott-exponential} and Elliott's absolute factor $12$,
this proves \eqref{eq:elliott-tv}.
\end{proof}

\subsection{The common good set}
The sector and total-factor statistics in \eqref{eq:statistics} and
\eqref{eq:wide-definition} are functions of this one divisor vector.
Their model means satisfy \eqref{eq:mean-bounds} and
\eqref{eq:wide-mean}; every failed threshold forces a deviation of
at least $t_0=L/16$. Lemma~\ref{lem:moment} and
\eqref{eq:normalized-moment} bound each model failure probability
by $4\,2^{-4r}$.

For each fixed $a$, take the union of all its failed sector tests
and its two failed total tests as a single event of the vector in
\eqref{eq:elliott-tv}. The comparison error is paid once for this
union, not once per character. Hence the enlarged bad set
$\mathcal B_N^*$ satisfies
\begin{align}
 \frac{|\mathcal B_N^*|}{H}
 &\le\sum_{a\in\mathcal A_N}
     \left(4(\varphi(a)+1)2^{-4r}
                +O_B((\log N)^{-B})\right)\notag\\
 &\ll_{B,\theta}
   (\log N)^{2\delta-2c\log2}
     +(\log N)^{\delta-B}
 \ll_\theta(\log N)^{-\delta},\label{eq:elliott-common-bad}
\end{align}
where $B>2\delta$ is fixed and $2c\log2>3\delta$.
Omitting the wider total test gives Proposition~\ref{prop:commonbad}.
Including it gives Proposition~\ref{prop:wide}, and the deterministic
argument of Section~4 completes the same local-count and multiplicity
conclusions. For each $a$ this is a joint approximation of its divisor
vector; shifts are combined by a union bound.

\section{The distribution inputs in the required normalization}\label{app:inputs}
We spell out the standard consequences used in Section~\ref{sec:good}. Write
$\Lambda_{\mathrm{vM}}$ for the von Mangoldt function, and let
\[
 \psi(t;\chi)=\sum_{n\le t}\Lambda_{\mathrm{vM}}(n)\chi(n),
 \qquad
 \vartheta(t;a,b)=\sum_{\substack{q\le t\\q\equiv b\ (a)}}\log q.
\]
For a reduced class, $\psi(t;a,b)$ denotes the analogous sum of
$\Lambda_{\mathrm{vM}}(n)$ over all integers in that class.

\subsection{Siegel--Walfisz and the harmonic means}
The character version on p.~2 of \cite{VaughanBV} gives, for fixed
$B>0$ and every character modulo $a\le(\log t)^B$,
\[
 \psi(t;\chi)-\one_{\chi=\chi_0}t
     \ll_B t\exp(-c_B\sqrt{\log t})
\]
with some $c_B>0$. Character orthogonality, equation~(6) there,
gives the same error for a reduced progression: the number of characters
cancels the factor $1/\varphi(a)$.\footnote{We use the character
statement; the preceding progression display in that version of the notes
omits the factor $t$ in its error term.}

The prime powers contribute $O(\sqrt t\log^2t)$, uniformly in the
progression. Partial summation (equation~(4), p.~1 of \cite{VaughanBV})
then gives \eqref{eq:sw}. For completeness, when $a\le\log t$, split
the integral at $\sqrt t$. In the upper part, $a\le(\log u)^2$ for
all $u\ge\sqrt t$ once $t$ is large, so the character bound applies
with the single fixed choice $B=2$. In the lower part the trivial bound
contributes at most $O(\sqrt t\log^2t)$. Exponential decay absorbs
every prescribed fixed power $(\log t)^{-H}$. The usual possible
ineffectivity of Siegel--Walfisz remains; it does not create an unproved
input or an effective onset.

In our application $a\le2(\log N)^\delta\le\log Q$ eventually.
Applying partial summation to $1/q$ on $(Q,Z]$ with $H=9$ bounds the
integrated error by $O((\log Q)^{-8})$. Replacing $1/q$ by
$1/(q-1)$ costs at most $O(\sum_{n>Q}n^{-2})=O(Q^{-1})$.
This proves \eqref{eq:lambda} with constants independent of $a$, its
unit class, and the growing order $r$.

\subsection{The dyadic prime-counting form of Bombieri--Vinogradov}
The two-term bound on p.~3 of \cite{VaughanBV} states, for fixed $A_0>0$,
\begin{align*}
 &\sum_{m\le M}\max_{(v,m)=1}\sup_{1\le u\le T}
   \left|\psi(u;m,v)-\frac{u}{\varphi(m)}\right|\\
 &\hspace{2cm}\ll_{A_0}
    T(\log T)^{-A_0}+T^{1/2}M\bigl(\log(TM)\bigr)^6.
\end{align*}
Use $T=2N$, $M=4N^{1/4}$, and denote the maximum and supremum
for a fixed $m$ by $E_m$. The second term is
$O(N^{3/4}(\log N)^6)$ and is smaller than every fixed logarithmic
saving relative to $N$.

For $N\le u\le2N$, put
$\Delta_m(u;v)=\vartheta(u;m,v)-u/\varphi(m)$. Removing prime powers
gives
\[
 \max_{(v,m)=1}\sup_{N\le u\le2N}|\Delta_m(u;v)|
       \le E_m+O(\sqrt N\log^2N).
\]
The exact dyadic partial-summation identity is
\begin{align*}
 &\pi(2N;m,v)-\pi(N;m,v)
       -\frac{\Li(2N)-\Li(N)}{\varphi(m)}\\
 &\quad=\frac{\Delta_m(2N;v)}{\log(2N)}
       -\frac{\Delta_m(N;v)}{\log N}
       +\int_N^{2N}\frac{\Delta_m(u;v)}{u\log^2u}\,du.
\end{align*}
After maximizing and summing, its absolute value is bounded by
\[
 \ll\frac1{\log N}\left(\sum_{m\le M}E_m
                           +M\sqrt N\log^2N\right).
\]
Choose $A_0=A+1$. Both terms of size $N^{3/4}$ are absorbed, giving
\eqref{eq:bv}. The proof uses only the interval $[N,2N]$ at the
partial-summation step, so no modulus-level assumption at much smaller
integration variables is needed.

\begin{thebibliography}{99}
\small
\raggedright
\bibitem{BE}
T. D. Browning and C. Elsholtz,
\emph{The number of representations of rationals as a sum of unit fractions},
Illinois J. Math. \textbf{55} (2011), 685--696.
\href{https://www.math.tugraz.at/~elsholtz/WWW/papers/papers33FINAL2013.pdf}{Author-hosted published version}.

\bibitem{EP}
C. Elsholtz and S. Planitzer,
\emph{The number of solutions of the Erd\H{o}s--Straus equation and sums
of $k$ unit fractions}, Proc. Roy. Soc. Edinburgh Sect. A
\textbf{150} (2020), 1401--1427.
\href{https://doi.org/10.1017/prm.2018.137}{doi};
\href{https://arxiv.org/abs/1805.02945v1}{arXiv:1805.02945v1}.

\bibitem{Chamberland}
M. Chamberland,
\emph{The Erd\H{o}s--Straus conjecture and the structure of primes},
Integers \textbf{26} (2026), Paper A42, 7 pp.
\href{https://math.colgate.edu/~integers/aa42/aa42.pdf}{Journal PDF}.

\bibitem{Bradford}
K. Bradford,
\emph{Elementary patterns from the Erd\H{o}s--Straus conjecture},
Integers \textbf{25} (2025), Paper A54, 10 pp.
\href{https://doi.org/10.5281/zenodo.15755958}{doi:10.5281/zenodo.15755958}.

\bibitem{Dahan}
B. Dahan,
\emph{Sieve dimension and search depth for the Erd\H{o}s--Straus conjecture,
$n\equiv1\pmod{24}$}, preprint, 2026,
\href{https://arxiv.org/abs/2608.24035v1}{arXiv:2608.24035v1}.

\bibitem{ET}
C. Elsholtz and T. Tao,
\emph{Counting the number of solutions to the Erd\H{o}s--Straus equation
on unit fractions}, J. Aust. Math. Soc. \textbf{94} (2013), 50--105.
Revised author version:
\href{https://arxiv.org/abs/1107.1010v6}{arXiv:1107.1010v6} (2015).

\bibitem{Bettin}
S. Bettin, J. Bober, A. Booker, B. Conrey, M. Lee, G. Molteni,
T. Oliver, D. J. Platt, and R. Steiner,
\emph{A conjectural extension of Hecke's converse theorem},
Ramanujan J. \textbf{47} (2018), 659--684, Lemma~4.9.
\href{https://doi.org/10.1007/s11139-017-9953-y}{doi};
\href{https://arxiv.org/abs/1704.02570}{arXiv:1704.02570}.

\bibitem{Hall}
R. R. Hall, \emph{The divisors of $p-1$},
Mathematika \textbf{20}(1) (1973), 87--97.
\href{https://doi.org/10.1112/S0025579300003661}{doi}.

\bibitem{Ford}
K. Ford,
\emph{Poisson approximation of prime divisors of shifted primes},
Int. Math. Res. Not. IMRN 2025(7), rnaf079.
\href{https://doi.org/10.1093/imrn/rnaf079}{doi};
\href{https://arxiv.org/abs/2408.03803v4}{arXiv:2408.03803v4}.

\bibitem{Gorodetsky}
O. Gorodetsky,
\emph{A Kubilius model for sieve-theoretic sequences},
Analysis Mathematica (2026), doi:10.1007/s10476-026-00178-w.
\href{https://doi.org/10.1007/s10476-026-00178-w}{doi};
\href{https://arxiv.org/abs/2608.14190v1}{arXiv:2608.14190v1}.

\bibitem{PW}
C. Pomerance and A. Weingartner,
\emph{Exceptions to the Erd\H{o}s--Straus--Schinzel conjecture},
Ramanujan J. \textbf{69} (2026), 31.
\href{https://doi.org/10.1007/s11139-025-01312-2}{doi};
\href{https://arxiv.org/abs/2511.16817}{arXiv:2511.16817}.

\bibitem{Bello}
M. Bello--Hern\'andez, M. Benito, and E. Fern\'andez,
\emph{A divisor parametrization for the Erd\H{o}s--Straus conjecture},
arXiv:2606.10922 (2026).
\href{https://arxiv.org/abs/2606.10922}{arXiv:2606.10922}.

\bibitem{VaughanBV}
R. C. Vaughan,
\emph{The Bombieri--Vinogradov theorem}, author lecture notes, Section~1.
\url{https://personal.science.psu.edu/rcv4/Bombieri.pdf}.
Accessed October 28, 2026.
\end{thebibliography}
\end{document}
